% Paper 4, §9 — Activation patching: the causal-intervention spine
% Thesis: activation patching is the canonical *causal* intervention
% method in mechanistic interpretability. The three-run protocol
% (clean / corrupted / patched) plus the indirect-effect statistic give
% the field its only direct causal handle on internal computation. Every
% subsequent intervention method (path patching, attribution patching,
% ACDC, attribution graphs) is a refinement of this protocol along one
% of three methodological axes — corruption method, metric, and
% granularity — that Zhang & Nanda 2023 systematize.

\section{Activation patching}
\label{sec:patching}

The lens and \glsterm{probe-2}{probe} of \S\ref{sec:logit-lens-family}--\S\ref{sec:probing}
read the \glsterm{residual-stream}{residual stream}; they tell us what intermediate states
\emph{contain}. They do not tell us what the model \emph{uses}. A
high-accuracy \glsterm{probe-2}{probe} at layer $\ell$ is consistent with the property
being computed at $\ell$ and consumed by the unembed, with the property
being computed at $\ell$ and ignored by the unembed (\emph{vestigial
information}), or with a confound entirely outside the model
\citep{belinkov2022-probing-survey}. Closing that gap requires an
\emph{intervention}: change the activation at $\ell$, hold everything
else fixed, observe the output. \glsterm{activation-patching}{Activation patching} is the canonical
form of that intervention. The three-run protocol of
\citet{meng2022-rome}, generalized and named by
\citet{zhang2023-apatching}, gives the field its workhorse causal
estimator and the heatmap that ROME made famous; the variant zoo of
\glsterm{path-patching}{path patching} \citep{wang2022-ioi}, \glsterm{attribution-patching}{attribution patching}, and ACDC
\citep{conmy2023-acdc} are all refinements of the same protocol along
three methodological axes that this section names explicitly.

\subsection{The three-run protocol}
\label{sec:patching-protocol}

\citet{meng2022-rome} fix the protocol as three forward
passes\footnote{KB excerpt:
\texttt{kb/excerpts/meng2022-rome.md\#sec-2}.}, restated by
\citet[\S 2.1]{zhang2023-apatching} as the now-standard
recipe\footnote{KB excerpt:
\texttt{kb/excerpts/zhang2023-apatching.md\#sec-2-1}.}. Fix a clean
prompt $X_{\text{cl}}$ with target answer $r$ (e.g.,
``\textit{The Eiffel Tower is in}'' $\to$ ``\textit{Paris}'') and a
corrupted prompt $X_{*}$ that changes the answer (e.g.,
``\textit{The Colosseum is in}'' $\to$ ``\textit{Rome}''). Let
$\mathbf{h}_i^{(\ell)}$ denote the model's residual-stream activation
at sequence position $i$ and layer $\ell$, and $C$ a candidate
\emph{component} indexing some collection of these activations
(a single \glsterm{hidden-state}{hidden state}, an MLP output, an attention-head output,
or a subspace; see \S\ref{sec:patching-granularity}).

\begin{enumerate}
  \item \textbf{Clean run.} Run $\mathcal{M}$ on $X_{\text{cl}}$ and
    cache every internal activation
    $\{\mathbf{h}_i^{(\ell)}\}_{i,\ell}$. Record the model's output
    distribution $\mathbb{P}_{\text{cl}}$.
  \item \textbf{Corrupted run.} Run $\mathcal{M}$ on $X_{*}$ and record
    the corrupted output distribution $\mathbb{P}_{*}$ and a target
    metric $m_{*}$ (probability, \glsterm{logit-difference}{logit difference}, or KL --- see
    \S\ref{sec:patching-metric}).
  \item \textbf{Patched run.} Run $\mathcal{M}$ on $X_{*}$, but at
    component $C$ overwrite the activation with the cached clean \glsterm{value}{value}
    from step~(1). Record the patched output distribution
    $\mathbb{P}_{\text{pt}}$ and the metric $m_{\text{pt}}$.
\end{enumerate}

The \emph{\glsterm{indirect-effect}{indirect effect}} of component $C$ on this prompt pair is the
metric difference between the patched and corrupted runs
\citep[\S 2]{meng2022-rome}\footnote{KB excerpt:
\texttt{kb/excerpts/meng2022-rome.md\#sec-2-effects}.}:
\begin{equation}
  \mathrm{IE}_C(X_{\text{cl}}, X_{*}) \;:=\; m_{\text{pt}} \;-\; m_{*}.
  \label{eq:ie-definition}
\end{equation}
Averaging over a dataset $\{(X_{\text{cl}}^{(n)}, X_{*}^{(n)})\}_{n=1}^{N}$
of paired counterfactuals gives the \emph{\glsterm{average-indirect-effect}{average indirect effect}}
\citep[\S 2]{meng2022-rome},
\begin{equation}
  \mathrm{AIE}_C \;:=\;
  \frac{1}{N} \sum_{n=1}^{N} \mathrm{IE}_C(X_{\text{cl}}^{(n)}, X_{*}^{(n)}),
  \label{eq:aie}
\end{equation}
which is the statistic the field actually reports: a single
$\mathrm{IE}_C$ on one prompt is noisy, $\mathrm{AIE}_C$ on a
counterfactual benchmark is the population-level localization signal.
A large $\mathrm{AIE}_C$ means substituting $C$'s clean activation back
into the corrupted run \emph{recovers} the clean behavior, so $C$ is a
sufficient causal mediator for the metric on this task.

Iterating $C$ over every (component, token, layer) cell of the
residual-stream grid produces the \emph{causal trace heatmap} of
\citet[Fig.~1]{meng2022-rome}: rows are token positions, columns are
layers, cells are $\mathrm{AIE}_C$. Bright cells localize the behavior;
dark cells are causally inert. The heatmap is the visual signature of
patching, and the substantive output of \S\ref{sec:rome} is what its
peaks revealed for factual recall.

\paragraph{Tensor-shape ledger.} A residual-stream activation at one
$(B, S)$ position is $\mathbf{h}^{(\ell, t)} \in \mathbb{R}^{D}$; the
batched tensor is $\BSH$. A
component $C$ identifies a subset of this grid: ``MLP output at
$(\ell, t)$'' is one $\mathbb{R}^{D}$ vector; ``attention-head $h$
output at $(\ell, t)$'' is one $\mathbb{R}^{d_h}$ vector before the
$W^O$ projection; ``head $h$ contribution to the \glsterm{residual-stream}{residual stream}'' is
the $\mathbb{R}^{D}$ vector after $W^O$. A subspace patch replaces only
the projection
$\mathbf{h}^{(\ell, t)} \mapsto \mathbf{h}^{(\ell, t)} -
(\mathbf{h}^{(\ell, t)})^{\top}\hat{\mathbf{d}}\;\hat{\mathbf{d}} +
(\mathbf{h}_{\text{cl}}^{(\ell, t)})^{\top}\hat{\mathbf{d}}\;
\hat{\mathbf{d}}$ along a \glsterm{probe-2}{probe} direction $\hat{\mathbf{d}}$, leaving
the orthogonal complement unchanged --- the natural granularity for
``patch the \glsterm{truth-direction-2}{truth direction}'' or ``patch the Paris \glsterm{feature}{feature}.''

\subsection{The metric: logit difference is the default}
\label{sec:patching-metric}

The metric $m$ in Eq.~\eqref{eq:ie-definition} is a free choice with
real consequences. \citet[\S 2.1]{zhang2023-apatching} catalog three
options\footnote{KB excerpt:
\texttt{kb/excerpts/zhang2023-apatching.md\#sec-2-1-metrics}.}.
Probability, $m^{\mathrm{P}} = \mathbb{P}(r)$, gives a unitless
patching effect $\mathbb{P}_{\text{pt}}(r) - \mathbb{P}_{*}(r)$; it is
the metric of \citet{meng2022-rome}. \glsterm{logit-difference}{Logit difference}, with $r'$ the
corrupted-prompt answer, is
\begin{equation}
  \mathrm{LD}(r, r') \;:=\; \mathrm{logit}(r) - \mathrm{logit}(r'),
  \label{eq:logit-diff}
\end{equation}
and the canonical normalized patching effect is
\citep[\S 2.1]{zhang2023-apatching}
\begin{equation}
  m^{\mathrm{LD}}_{\text{pt}} \;:=\;
  \frac{\mathrm{LD}_{\text{pt}}(r, r') - \mathrm{LD}_{*}(r, r')}{
        \mathrm{LD}_{\text{cl}}(r, r') - \mathrm{LD}_{*}(r, r')},
  \label{eq:logit-diff-normalized}
\end{equation}
typically in $[0, 1]$ (where $1$ is fully restored, $0$ is corrupted).
KL is
$D_{\mathrm{KL}}(\mathbb{P}_{\text{cl}} \,\|\, \mathbb{P}_{*}) -
 D_{\mathrm{KL}}(\mathbb{P}_{\text{cl}} \,\|\, \mathbb{P}_{\text{pt}})$.

The methodological case for \glsterm{logit-difference}{logit difference} as the default is made by
\citet[\S 4.1]{heimersheim2024-apatching-howto}\footnote{KB excerpt:
\texttt{kb/excerpts/heimersheim2024.md\#sec-4-1}.}: it is a
``mostly linear function of the \glsterm{residual-stream}{residual stream} (unlike
probability-based metrics) which makes it easy to directly attribute
\glsterm{logit-difference}{logit difference} to individual components,'' and it is ``a `softer'
metric, allowing us to see partial effects on the model even if they
don't change the rank of the output tokens.''
\citet[\S 4.2]{heimersheim2024-apatching-howto} catalog the failure
modes of the alternatives\footnote{KB excerpt:
\texttt{kb/excerpts/heimersheim2024.md\#sec-4-2}.}: probability is
``non-linear, in the sense that [it] track[s] the \glsterm{logits}{logits} exponentially.
For example, a model component adding $+2$ to a given logit can create
a $1$ or $40$ percentage point probability increase, just depending on
what the baseline was''; log-probabilities saturate (``once the correct
answer becomes the model's top guess, the logprob stops increasing
meaningfully''); rank- and accuracy-based metrics ``round off each
input to a discrete metric.'' The single most important consequence of
choosing probability over \glsterm{logit-difference}{logit difference} is invisibility of negative
components: the IOI Negative Name Mover heads
\citep[\S 3]{wang2022-ioi}\footnote{KB excerpt:
\texttt{kb/excerpts/wang2022-ioi.md\#sec-3-\glsterm{circuit}{circuit}}.} push \emph{against}
the correct answer, and a probability metric near $\mathbb{P}(r) \in
\{0, 1\}$ saturates exactly the regime in which their effect is
detectable. \glsterm{logit-difference}{Logit difference} does not.

\subsection{The corruption: STR over GN}
\label{sec:patching-corruption}

The second methodological axis is how $X_{*}$ is generated from
$X_{\text{cl}}$. \citet[\S 2.1]{zhang2023-apatching} fix two
canonical choices\footnote{KB excerpt:
\texttt{kb/excerpts/zhang2023-apatching.md\#sec-2-1-corruption}.}.
\textbf{\glsterm{gaussian-noising}{Gaussian Noising}} (GN), used by \citet{meng2022-rome}, ``adds
Gaussian noise $\mathcal{N}(0, \nu)$ to the embeddings of certain \glsterm{key}{key}
tokens, where $\nu$ is 3 times the standard deviation of the token
embeddings from the textset.'' \textbf{Symmetric Token Replacement}
(STR), used by \citet{wang2022-ioi}, ``replaces the \glsterm{key}{key} tokens by
similar ones with equal sequence length''; ``Eiffel Tower'' becomes
``Colosseum'', ``Mary'' becomes ``Alice''. STR's defining property is
that ``$X_{\mathrm{corrupt}}$ is identically distributed as a fresh
draw of a clean prompt'' \citep[\S 2.1]{zhang2023-apatching}.

\begin{contradiction}
\textbf{GN can mis-localize relative to STR.}
\citet[\S 1]{zhang2023-apatching} report that ``GN and STR can lead to
inconsistent localization and \glsterm{circuit}{circuit} discovery outcomes (Section
3.1)'' on the same task. The conjectured mechanism: ``GN breaks
model's internal mechanisms by putting it off
distribution''\footnote{KB excerpt:
\texttt{kb/excerpts/zhang2023-apatching.md\#sec-findings}.} ---
in-distribution prompts that the model knows how to process produce
clean signals about which components matter; out-of-distribution noise
produces signals about which components fail under noise, which need
not coincide. \citet[\S 1]{zhang2023-apatching} ``advocate for STR, as
it supplies in-distribution corrupted prompts that help to preserve
consistent model behavior.'' Concrete consequence:
\citet{meng2022-rome}'s GN-derived \glsterm{average-indirect-effect}{AIE} peak at GPT-2-XL layer $15$
broadly survives STR re-validation \citep[\S 3]{zhang2023-apatching},
but the precise per-layer peak shifts, and downstream uses of ROME's
exact layer-of-edit recommendation inherit that sensitivity. The
strongest form of the methodology claim --- ``GN-derived localizations
should be re-validated under STR before being treated as canonical
for a new task'' --- is the field's current default.
\end{contradiction}

A practical refinement, due to
\citet[\S 2.3]{heimersheim2024-apatching-howto}, is the
\emph{denoising vs.\ noising} distinction\footnote{KB excerpt:
\texttt{kb/excerpts/heimersheim2024.md\#sec-2-3}.}. The protocol
above is \emph{denoising} (clean $\to$ corrupted): the patched run
inserts a clean activation into a corrupted run and asks
\emph{which patches were sufficient to restore the behavior}. The
inverse protocol, \emph{noising} (corrupted $\to$ clean), inserts a
corrupted activation into a clean run and asks
\emph{which patches were necessary to maintain the behavior}.
\citet[\S 2.4]{heimersheim2024-apatching-howto} show with a worked
AND/OR-gate example that the two directions are not symmetric: an AND
\glsterm{circuit}{circuit} ($C = A \wedge B$) has both $A$ and $B$ as necessary, so
noising flags both; denoising of either alone leaves the other at the
corrupted baseline and finds nothing. An OR \glsterm{circuit}{circuit} ($C = A \vee B$)
has either alone as sufficient, so denoising of either restores
behavior; noising of either alone leaves the other operative. The
practical recommendation is to choose direction by the suspected
\glsterm{circuit}{circuit} topology, or to run both.

\subsection{Granularity: from hidden states to subspaces}
\label{sec:patching-granularity}

The third axis is what counts as $C$. The granularities used in the
literature \citep[\S 2.1]{zhang2023-apatching} form a coarse-to-fine
ladder, summarized in Table~\ref{tab:patching-granularity}.

\begin{table}[h]
\centering
\small
\begin{tabular}{@{}lll@{}}
\toprule
Granularity & Component $C$ & Used by \\
\midrule
\glsterm{hidden-state}{Hidden state} & $\mathbf{h}_i^{(\ell)} \in \mathbb{R}^{D}$ & \citet{meng2022-rome} \\
MLP output & $\mathrm{MLP}^{(\ell)}(\mathbf{h}_i^{(\ell)}) \in \mathbb{R}^{D}$ & \citet{meng2022-rome} \\
Attention-layer output & $\mathrm{Attn}^{(\ell)}(\mathbf{h}_i^{(\ell)}) \in \mathbb{R}^{D}$ & \citet{meng2022-rome} \\
Single attention head & $h_{j}^{(\ell)} \in \mathbb{R}^{d_h}$ before $W^{O}$ & \citet{wang2022-ioi} \\
Subspace & $\mathbf{P}_{\hat{\mathbf{d}}}\,\mathbf{h}_i^{(\ell)}$ & \glsterm{probe-2}{probe}-direction patches \\
\bottomrule
\end{tabular}
\caption{Patching granularities. \citet{meng2022-rome} patch hidden
states, MLPs, and attention layers separately to attribute factual
recall to MLPs over attention; \citet{wang2022-ioi} patch individual
heads to recover the seven-class IOI circuit; subspace patching
along a probe direction $\hat{\mathbf{d}}$ tests whether a specific
\emph{property} encoded by a probe is causally used.}
\label{tab:patching-granularity}
\end{table}

The headline trade-off: granularity governs the strength of the claim
the experiment can support. Hidden-state patching localizes
\emph{layers} and \emph{token positions}; MLP-vs-attention patching
splits the contribution by sublayer; head-level patching localizes
specific attention heads and is the granularity needed for \glsterm{circuit}{circuit}
claims; subspace patching localizes specific representational
directions. Coarse patching that finds ``MLPs at layer 15 matter''
does not yet justify ``the \texttt{Paris} direction in those MLPs
matters,'' even if both are true --- the latter requires either
finer-grained patching at the subspace level or a downstream \glsterm{probe-2}{probe}-on-%
patched-states design \citep[\S 2.3]{zhang2023-apatching}.

\subsection{Path patching: isolating direct effects}
\label{sec:patching-path}

Standard \glsterm{activation-patching}{activation patching} gives the \emph{total} effect of a
component on the metric: any influence routed through $C$ is captured,
including effects mediated by downstream components. Circuits with
multiple parallel paths --- which is most circuits --- need a finer
estimate. \citet[\S 3.1]{wang2022-ioi}'s \emph{path
patching}\footnote{KB excerpt:
\texttt{kb/excerpts/wang2022-ioi.md\#sec-3-1-path-patching}.} is the
partial-derivative analogue: it isolates the direct effect of a head
$h$ on a downstream component $R$ by patching only along the residual
and MLP paths from $h$ to $R$, holding indirect routes through other
attention heads at their corrupted values. Verbatim from
\citet[\S 3.1]{wang2022-ioi}:

\begin{quote}
Given inputs $x_{\text{orig}}$ and $x_{\text{new}}$, and a set of paths
$\mathcal{P}$ emanating from a node $h$, \glsterm{path-patching}{path patching} runs a forward
pass on $x_{\text{orig}}$, but for the paths in $\mathcal{P}$ it
replaces the activations for $h$ with those from $x_{\text{new}}$. In
our case, $h$ will be a fixed attention head and $\mathcal{P}$
consists of all direct paths from $h$ to a set of components $R$, i.e.\
paths through residual connections and MLPs (but not through other
attention heads); this measures the counterfactual effect of $h$ on
the members of $R$.
\end{quote}

The methodological gain is exactly what makes the seven-class IOI
decomposition (\S\ref{sec:circuits}) legible: full \glsterm{activation-patching}{activation patching}
attributes ``Name Mover Heads matter,'' but only \glsterm{path-patching}{path patching} can
attribute ``S-Inhibition Heads write into the queries of the Name
Mover Heads,'' i.e., a directed effect $h \to R$ via a specific
projection. Without \glsterm{path-patching}{path patching} the IOI \glsterm{circuit}{circuit} collapses into a
flat ranked list of important heads; with it, it is a wiring diagram.

\begin{intuition}
\glsterm{activation-patching}{Activation patching} is counterfactual surgery on the residual-stream
graph. The clean run is the patient who knows the answer; the
corrupted run is the patient who lost it; the patched run grafts one
specific tissue from the clean patient into the corrupted patient. If
the corrupted patient recovers, that tissue carried the answer.
Returning to canonical math: the recovery measure is exactly
$\mathrm{IE}_C = m_{\text{pt}} - m_{*}$ from
Eq.~\eqref{eq:ie-definition}, with $C$ the grafted tissue and $m$ a
behaviorally meaningful metric (\glsterm{logit-difference}{logit difference}, by default). Path
patching extends the surgery to grafting along specific
\emph{conduits} between tissues rather than the tissues themselves;
the canonical math is the same Eq.~\eqref{eq:ie-definition} restricted
to a path set $\mathcal{P}$.
\end{intuition}

\subsection{Attribution patching: the linearization}
\label{sec:patching-attribution}

The cost of full \glsterm{activation-patching}{activation patching} is one forward pass per component:
sweeping every (component, token, layer) cell of the heatmap costs
$O(N_{\text{components}})$ forward passes per (clean, corrupted) pair,
and \glsterm{path-patching}{path patching} at depth-$k$ costs $O(N_{\text{components}}^{k})$.
On GPT-2 small this is feasible; on production-scale models it is not.
\emph{\glsterm{attribution-patching}{Attribution patching}} replaces the substitution by its first-order
Taylor approximation: write
$\Delta\mathbf{h}_C := \mathbf{h}_{C,\text{cl}} -
\mathbf{h}_{C,*}$ for the clean-minus-corrupted activation difference at
component $C$, and approximate the \glsterm{indirect-effect}{indirect effect} of patching $C$ as
\begin{equation}
  \mathrm{IE}_C
  \;\approx\;
  \big\langle \nabla_{\mathbf{h}_C} m \,\big|_{\mathbf{h}_C =
  \mathbf{h}_{C,*}},\;
  \Delta\mathbf{h}_C \big\rangle,
  \label{eq:attribution-patching}
\end{equation}
where the gradient is taken at the corrupted activation. The cost is
one forward and one backward pass on the corrupted prompt to obtain
$\nabla_{\mathbf{h}} m$ at \emph{every} component simultaneously, plus
one cached clean forward pass for $\Delta\mathbf{h}$; this collapses
$O(N_{\text{components}})$ patch-forward-passes into $O(1)$
forward+backward passes, an asymptotic win that scales the technique
to production models.

\begin{forumsignal}
The attribution-patching approximation in
Eq.~\eqref{eq:attribution-patching} was popularized by Nanda's
``\glsterm{attribution-patching}{Attribution Patching}: \glsterm{activation-patching}{Activation Patching} at Industrial Scale''
LessWrong post (2023), which is the canonical reference for the
technique \deepencite{nanda2023-attribution-patching:
LessWrong / Tier B; awaiting peer-reviewed citation}. Subsequent
peer-reviewed venues (notably \citealt{conmy2023-acdc} and the
Anthropic attribution-graph line of \S\ref{sec:circuits}) build on
the linearization, but the original write-up is a forum post and
should be flagged as such. Hard claims about attribution-patching's
\emph{accuracy} relative to full patching back to
\citet[\S 1]{zhang2023-apatching}, who note the approximation
``is known to fail in non-linear regions''\footnote{KB excerpt:
\texttt{kb/excerpts/zhang2023-apatching.md\#sec-1}.} ---
attention-\glsterm{softmax}{softmax} \glsterm{saturation}{saturation} and \glsterm{relu}{ReLU}/GeLU dead zones being the
canonical failure cases.
\end{forumsignal}

The mediation-analysis lineage of \glsterm{activation-patching}{activation patching}, predating its
mech-interp adoption, traces to causal-mediation work in
NLP \deepencite{vig2020-causal-mediation §3:
awaiting peer-reviewed bibtex entry}. \citet[\S 1]{meng2022-rome} cite
this lineage explicitly when introducing the three-run causal-tracing
protocol; \citet[\S 1]{zhang2023-apatching} list ``causal mediation
analysis'' as one of the synonyms under which \glsterm{activation-patching}{activation patching}
appears in the literature. The genealogy matters because it locates
patching in a pre-existing statistical-causal \glsterm{vocabulary}{vocabulary}
(total/\glsterm{indirect-effect}{indirect effect}, mediation, intervention) that the
mech-interp community then specialized to neural-network
computational graphs.

\subsection{Recommendations and pitfalls}
\label{sec:patching-recommendations}

The methodology of \citet{zhang2023-apatching} and
\citet{heimersheim2024-apatching-howto} converges on a default recipe.
Table~\ref{tab:patching-recommendations} summarizes.

\begin{table}[h]
\centering
\small
\begin{tabular}{@{}lll@{}}
\toprule
Choice & Default & Source \\
\midrule
Corruption & STR (in-distribution) & \citet[\S 1]{zhang2023-apatching} \\
Direction & Denoising; noising for AND-circuits & \citet[\S 2.4]{heimersheim2024-apatching-howto} \\
Metric & \glsterm{logit-difference}{Logit difference} (normalized) & \citet[\S 4.1]{heimersheim2024-apatching-howto} \\
Granularity & Head-level for circuits; subspace for features & \citet[\S 2.3]{zhang2023-apatching} \\
Validation & Pair patching with a complementary ablation test & \citet[\S 4]{wang2022-ioi} \\
\bottomrule
\end{tabular}
\caption{Default activation-patching recipe consolidated from the
methodology literature. Each row has at least one peer-reviewed
recommendation; the validation row reflects the
faithfulness/completeness/minimality discipline of
\citet[\S 4]{wang2022-ioi}, which patching alone does not enforce.}
\label{tab:patching-recommendations}
\end{table}

Three pitfalls deserve naming, lifted from
\citet[\S 4.2]{heimersheim2024-apatching-howto} and
\citet[\S 2.1]{zhang2023-apatching}. \emph{Logit-difference
false-positives}: because \glsterm{logit-difference}{logit difference} is a difference, ``it may
be driven by either getting better at the correct answer or worse at
the incorrect answer''\footnote{KB excerpt:
\texttt{kb/excerpts/heimersheim2024.md\#sec-4-2}.} --- a patch that
indiscriminately damages the model toward uniform output reduces the
incorrect-answer logit and inflates the apparent patching effect,
without localizing anything specific. The fix is sanity-checking
against an absolute-probability metric and against unrelated control
metrics, not abandoning \glsterm{logit-difference}{logit difference}. \emph{Probability
\glsterm{saturation}{saturation}}: when $\mathbb{P}(r) \in \{0, 1\}$ at either endpoint, the
patching effect on probability is artificially compressed; logit
difference is the explicit fix. \emph{Coarse-granularity sleights of
hand}: a heatmap showing \glsterm{average-indirect-effect}{AIE} concentrated at ``MLP layer 15, last
subject token'' is consistent with the model storing the fact in that
MLP, with the model storing the fact upstream and reading it at that
MLP, with the model storing the fact in a \glsterm{feature}{feature} subspace within
that MLP, and with the model storing the fact partly there and partly
in a redundant location whose \glsterm{average-indirect-effect}{AIE} is masked by backup
\citep[\S 3]{wang2022-ioi}. Each of these claims requires a
finer-grained or causally complementary experiment to distinguish.

\subsection{Forward link}
\label{sec:patching-forward}

\glsterm{activation-patching}{Activation patching} is the method that turns the lens-and-\glsterm{probe-2}{probe}
correlational reading of the \glsterm{residual-stream}{residual stream}
(\S\ref{sec:logit-lens-family}--\S\ref{sec:probing}) into causal claims
the field can adjudicate. The next two sections deploy it.
Section~\ref{sec:rome} follows the heatmap of
\citet[Fig.~1]{meng2022-rome} from a localization claim
(``mid-layer MLPs at the last subject token mediate factual recall'')
to a weight-level intervention (the rank-one MLP edit) that turns
patching's diagnostic output into capability surgery; this is the
bridge from interpretability-method to model-editing
methodology. Section~\ref{sec:circuits} consumes patching as one of two
methods that produce the IOI \glsterm{circuit}{circuit} and the ACDC subgraph: Wang et
al.'s manual path-patching reverse-engineering and Conmy et al.'s
KL-thresholded automated edge-pruning are both refinements of the
three-run protocol, and the cross-method evidence with \glsterm{sparse-autoencoder}{SAE} features
and the \glsterm{tuned-lens}{tuned lens} (\S\ref{sec:cross-method}) is what gives the IOI
picture its load-bearing status. \glsterm{activation-patching}{Activation patching} is, in this
sense, the spine: every causal interpretability claim downstream of
this section either uses it directly or refines it.
